\section{Traditional Venue Integration}
\label{sec:traditional}

\subsection{Alpaca OmniSub}

Alpaca serves as the primary traditional equity execution provider. The integration uses the Alpaca Broker API for account management and order routing:

\begin{lstlisting}[language=Go,caption={Alpaca provider implementation.}]
type AlpacaProvider struct {
    client *alpaca.Client
    mu     sync.RWMutex
}

func (a *AlpacaProvider) Submit(
    ctx context.Context, order Order,
) (*Execution, error) {
    req := alpaca.PlaceOrderRequest{
        Symbol:      order.Symbol,
        Qty:         decimal.NewFromInt(int64(order.Quantity)),
        Side:        mapSide(order.Side),
        Type:        mapOrderType(order.Type),
        TimeInForce: mapTIF(order.TimeInForce),
    }

    if order.Type == Limit {
        req.LimitPrice = &order.Price
    }

    alpacaOrder, err := a.client.PlaceOrder(req)
    if err != nil {
        return nil, fmt.Errorf("alpaca submit: %w", err)
    }

    return &Execution{
        OrderID:    order.ID,
        Venue:      a.ID(),
        ExternalID: alpacaOrder.ID,
        Status:     mapStatus(alpacaOrder.Status),
    }, nil
}

func (a *AlpacaProvider) SettlementTime() time.Duration {
    return 24 * time.Hour // T+1 for equities
}
\end{lstlisting}

\subsection{Settlement Reconciliation}

For orders executed on traditional venues, the SOR reconciles settlement between the off-chain execution and on-chain records:

\begin{enumerate}[nosep]
  \item Traditional venue executes the order (T+0 execution).
  \item SOR records execution in the \texttt{executions} collection.
  \item Traditional venue settles (T+1 for equities).
  \item SOR mints/burns tokenized wrapper on Lux EVM to reflect settlement.
  \item On-chain state reconciled with off-chain position.
\end{enumerate}

